%%%%%%%%%%%%%%%%%%%% author.tex %%%%%%%%%%%%%%%%%%%%%%%%%%%%%%%%%%%
%
% sample root file for your "contribution" to a proceedings volume
%
% Use this file as a template for your own input.
%
%%%%%%%%%%%%%%%% Springer %%%%%%%%%%%%%%%%%%%%%%%%%%%%%%%%%%


\documentclass{svproc}
%
% RECOMMENDED %%%%%%%%%%%%%%%%%%%%%%%%%%%%%%%%%%%%%%%%%%%%%%%%%%%
%

% to typeset URLs, URIs, and DOIs
\usepackage{url}
\def\UrlFont{\rmfamily}
\usepackage{cite}
\usepackage{cite}
\usepackage{amsmath,amssymb,amsfonts}
\usepackage{algorithmic}
\usepackage[final]{graphicx}
\usepackage{graphics}
\usepackage{textcomp}
\usepackage{xcolor}
\usepackage{tabularx}
\usepackage{hyperref}
\begin{document}
\mainmatter              % start of a contribution
%
\title{Artificial intelligence for enhancing the quality of
higher education in both online or offline: A
systematic Literature Review.}
%
\titlerunning{SLR}  % abbreviated title (for running head)
%                                     also used for the TOC unless
%                                     \toctitle is used
%
\author{Muwaffaq Imam\inst{1}}

% \author{Muwaffaq Imam\inst{1} \and Roger Temam\inst{2}
% Jeffrey Dean \and David Grove \and Craig Chambers \and Kim~B.~Bruce \and
% Elsa Bertino}
%
\authorrunning{Muwaffaq Imam} % abbreviated author list (for running head)
%
%%%% list of authors for the TOC (use if author list has to be modified)

\tocauthor{Ivar Ekeland, Roger Temam, Jeffrey Dean, David Grove,
Craig Chambers, Kim B. Bruce, and Elisa Bertino}
%
\institute{Innopolis University, Kazan, Russia,\\
\email{m.imam@innopolis.university},\\ WWW home page:
\texttt{https://innopolis.university/en/}
}

% \tocauthor{Ivar Ekeland, Roger Temam, Jeffrey Dean, David Grove,
% Craig Chambers, Kim B. Bruce, and Elisa Bertino}
% %
% \institute{Princeton University, Princeton NJ 08544, USA,\\
% \email{I.Ekeland@princeton.edu},\\ WWW home page:
% \texttt{http://users/\homedir iekeland/web/welcome.html}
% \and
% Universit\'{e} de Paris-Sud,
% Laboratoire d'Analyse Num\'{e}rique, B\^{a}timent 425,\\
% F-91405 Orsay Cedex, France}

\maketitle              % typeset the title of the contribution

\begin{abstract}
\label{Abstract}
\textbf{Purpose:} This Systematic Literature Review (SLR) aims to investigate the potential of Artificial Intelligence (AI) in enhancing the quality of education. Our focus is directed towards scrutinizing existing research within the field of AI in education, specifically considering the methodologies employed, the metrics utilized to measure educational quality post-implementation, and the impact on educational facilities.

\textbf{Methods:} To achieve our objectives, we strictly adhered to best practices in our field and conducted a comprehensive SLR. Our search methodology initially identified a promising collection of 2,829 potentially relevant papers. These were then reviewed according to pre-established inclusion and exclusion criteria, leading to the selection of 30 research papers that were incorporated into our final compilation.

\textbf{Results:} Our analysis resulted in the identification of five distinct use cases in which researchers employed AI to provide insights into enhancing educational quality. Additionally, we collected the metrics used by these researchers and categorized them into four broad categories. Furthermore, we gathered all the factors that researchers used to calculate these metrics, classifying these factors into four categories. We also examined the educational components that were influenced by these AI techniques, subsequently categorizing them into four groups based on their similarities. Lastly, it was observed that only three researchers, out of all the papers we reviewed, had developed tools intended for community use, further underscoring an area for potential future exploration and development.

% We would like to encourage you to list your keywords within
% the abstract section using the \keywords{...} command.
\keywords{Artificial intelligence, Education, SLR}
\end{abstract}
%

\section{Introduction}
\label{introduction}
Artificial Intelligence (AI) has emerged as a transformative force in various sectors, and education is no exception.

Our initial hypothesis suggests that AI can enhance education. In this study, we explore the methods through which AI can achieve this enhancement and the potential impact of these methods on educational facilities. These research questions guided our Systematic Literature Review (SLR).

To conduct the systematic literature review, we reviewed previous SLRs related to AI in education. According to the referenced research \cite{Artificial_intelligence_in_online_higher_education:_A_systematic_review_of_empirical}, educational institutions began exploring the integration of AI in education, particularly between 2019 and 2021. Numerous applications have been conducted within these institutions, with one of the most common being the prediction of student performance based on grades. This application can also extend to predicting dropout risks, forecasting student academic performance, and determining student satisfaction with online courses \cite{Artificial_intelligence_in_online_higher_education:_A_systematic_review_of_empirical} \cite{AI_IN_Education_SLR}.

Furthermore, research has categorized the usage of AI in educational facilities into four main categories \cite{Systematic_literature_review_on_opportunities_challenges_and_future_research_recommendations_of_artificial_intelligence_in_education}:

Firstly, in the context of learning, AI can be leveraged to evaluate student feedback and assignments, generating valuable insights to refine the learning process. Secondly, in the realm of teaching, AI enables the creation of adaptive instructional strategies, customizing the educational experience according to each student's unique needs and learning methods.

The third application is in assessment, where AI enables the integration of automatic grading systems, making the assessment process more efficient and providing prompt feedback to students. Lastly, AI plays a crucial role in administration, enhancing the efficiency of management systems in educational institutions and simplifying a wide range of administrative duties.

Machine Learning (ML) techniques have been applied with positive effects in higher education, particularly in predicting student performance, enhancing the quality of education, and reducing student attrition. ML algorithms are being used to assess students and provide continuous feedback to teachers, both in a classroom and blended learning setting. Several studies have explored ML-based approaches to combat cheating and ensure the academic integrity of student assessments \cite{AI_IN_Education_SLR}.

In their study \cite{Artificial_intelligence_in_online_higher_education:_A_systematic_review_of_empirical}, four major functions of AI applications in online higher education were identified:

Firstly, they enable predictive analysis, accurately forecasting student performance and satisfaction levels in both online and traditional classes. Secondly, they facilitate personalized resource recommendations. AI-based systems recommend relevant learning resources to online students, taking into account their preferences, interests, prior knowledge, and memory capacities.

Thirdly, automatic assessment is another key function. Researchers have developed automated essay evaluation systems, providing immediate assessment, feedback, and automated grading in online English learning platforms.

Finally, AI applications significantly enhance the learning experience. For instance, through the creation of Virtual Reality (VR) tools integrated with AI technology, especially in history learning. These VR tools allow students to fully engage with virtual environments, emulating real cities, where they can learn by exploring and interacting with virtual inhabitants.

The research conducted in this field has yielded promising results and has the potential to significantly impact educational facilities. These researchers have inspired our work in this SLR to delve deeper into enhancing the quality of education. The structure of this paper is as follows: We present our search protocol in Section \ref{Review_PROTOCOL}. Subsequently, we address our research questions in Section \ref{Results_and_Discussion}. A summary of our findings, highlighting the significance of this research, is provided in Section \ref{SYNOPTIC_SUMMARY_AND_SIGNIFICANCE_OF_RESEARCH}. Lastly, we discuss the limitations of our study, potential threats to validity, and the review assessment in Section \ref{LIMITATIONS,_THREATS_TO_VALIDITY,_AND_REVIEW_ASSESSMENT}.


\section{REVIEW PROTOCOL}
\label{Review PROTOCOL}
This study was conducted using a systematic literature review methodology, following the guidelines set forth by the PRISMA checklist. Additionally, we adhered to the research protocol outlined in \cite{Technologies_for_GQM_Based_Metrics_Recommender_Systems_A_Systematic_Literature_Review}. We initiated our study by formulating precise research questions that served as the foundation for our investigation.

Subsequently, in alignment with these questions, we crafted a comprehensive search query. This detailed query was employed across three distinct digital libraries, thereby ensuring access to a wide range of potential sources for inclusion.

Upon collecting the resources, we applied stringent inclusion and exclusion criteria to filter the literature. This rigorous screening process ensured that only the most relevant and high-quality papers were included for further analysis.

Finally, after meticulous evaluation and filtering, we arrived at a final selection of key papers for thorough reading and review. These articles formed the cornerstone of our study, providing in-depth insights that effectively addressed our research questions.

\subsection{PRISMA Checklist}
The full PRISMA checklist can be found in Appendix \ref{Appendix}.

\subsection{Research Questions}
Our research was based on the hypothesis: Can we use AI to enhance education, and to what extent can these applications impact our educational facilities? Based on this hypothesis, we formulated the following research questions:

\begin{enumerate}
\item[RQ1.] What are the use cases in which AI is employed to gain insights into enhancing education? Additionally, what AI algorithms and techniques are commonly used for this purpose? \label{RQ1}
\item[RQ2.] What are the metrics commonly utilized to assess the quality of both online and offline education? Furthermore, what factors are considered in the evaluation of these metrics?
\item[RQ3.] What is the impact of AI on the quality of education in both online and offline teaching formats?
\end{enumerate}


In this study, we aim to explore the role of artificial intelligence (AI) in enhancing education and its impact on educational quality in both online and offline teaching formats. Our first research question (RQ1) focuses on identifying the various use cases of AI that have been employed to gain insights and improve educational practices. We will investigate how AI algorithms and techniques are utilized in educational settings to enhance teaching and learning experiences. RQ2 examines the metrics commonly used to assess the quality of both online and offline education. We will investigate the factors considered in evaluating these metrics, such as student engagement, learning outcomes, instructor effectiveness, and technological infrastructure. Finally, RQ3 aims to understand the overall impact of AI on the quality of education in both online and offline formats. We will examine the effects of AI implementation on student performance, engagement, and satisfaction, as well as the potential benefits and challenges associated with AI integration in educational environments. Through this research, we seek to gain valuable insights into the use of AI in education and its implications for educational quality.


\subsection{Search process}
\label{searchStrategy}
 The databases we used include: Scopus, IEEE Xplore, and ACM Digital Library. W



\textbf{ Search Queries \& initial results:}

To ensure that the search results for measuring the quality of higher education courses using AI-related technology we divided the query into three sections 1)AI- technologies 2)Quality measures 3)higher education synonyms. For example, we can use search terms such as "artificial intelligence", "machine learning", "natural language processing", or "deep learning" for the AI technologies.
For the quality measures section, some possible search terms could be "course quality", "learner satisfaction", "teaching effectiveness", or "MOOC satisfaction".
To explore higher education in both online and offline contexts, we can use synonyms such as "post-secondary education", "tertiary education", "university education", or "college education".
Putting it all together, a potential search query format could be: "AI technology AND quality measures AND higher education synonym".

("Text mining" OR "sentiment analysis" OR "opinion mining" OR "NLP" OR "AI" OR "supervised machine learning algorithm" OR "sentiment analysis" OR "AIEd" OR "Artificial Intelligence in Education" OR "machine learning" OR "neural network" OR "deep learning" OR "NLP" OR "Natural language processing" OR "unsupervised machine learning algorithm" OR "text analysis" OR "Artificial Intelligence in Education")
AND
("course quality" OR "user satisfaction" OR "learner satisfaction" OR "student satisfaction" OR "teaching quality" OR "course satisfaction" OR "MOOCs satisfaction" OR "MOOC satisfaction")
AND
("online education" OR "online learning" OR "e-learning" OR MOOC* OR SPOC* OR "blended learning" OR "higher education" OR "online higher education" OR "undergraduate education" OR "graduate education" OR "university" OR "offline education" OR "postgraduate" OR "post-secondary education" OR "tertiary education" OR "higher education", "Online courses", "Online course" OR "college education")

\textbf{ScienceDirect Query: } 
since in scienceDirect we can't use advanced search, below is the query we used. 
"AIEd" OR "Artificial Intelligence in Education" OR "machine learning" OR "deep learning" OR "Natural language processing" OR "education" OR MOOCS
in abstract 
"course quality" OR "user satisfaction" OR "learner satisfaction" OR "student satisfaction" OR "teaching quality" OR "course satisfaction" OR "MOOCs satisfaction" OR "MOOC satisfaction"

WE haven't used another query because the results from the mentioned query produced many results. as we will discuss further in the up coming section.

In our Systematic Literature Review (SLR), we've amassed a substantial corpus of papers from three renowned digital libraries. A total of 424 papers were retrieved from ScienceDirect, marking the highest contribution. This was closely followed by 343 papers from the Association for Computing Machinery (ACM) database. Meanwhile, the IEEE Explore database yielded a robust selection as well, with 316 papers. These figures represent our comprehensive and far-reaching exploration of the available literature on the topic, indicating the thoroughness of our SLR. The comprehensive breakdown of these results can be found in the detailed SLR figure we've compiled. Everything is illustrated in figure \ref{SLR diagram}
\begin{enumerate}
    \item \textbf{ScienceDirect (424 papers)} 
    \item \textbf{IEEE Explore (316 papers)}
    \item \textbf{ACM (343 papers)}
\end{enumerate}

\begin{figure}[htbp]
\centerline{\includegraphics[scale=0.7]{figs/SLR diagram.jpeg}}
\caption{SLR diagram}
\label{SLR diagram}
\end{figure}













\subsection{Inclusion and exclusion criteria}
\label{IC_IE}
We formulated Exclusion Criteria (EC) and Inclusion Criteria (IC). EC and IC help deciding what research to include to our SLR.

\textbf{Inclusion Criteria (EC)} as follow: 
\begin{enumerate}

\item [IC1-] Research must report the application of AI in enhance or measure the quality of education. 
\item [IC2-] Research must reported as empirical research to demonstrate how they used AI algorithm in measuring course quality.
\item [IC3-] Research must report what AI algorithms/methods/system they used conducting the research.
\item [IC4-] Research must report the metrics they used to measure the quality of the course.

\end{enumerate}
\textbf{Exclusion Criteria (EC)} as follow: 
\begin{enumerate} 
\item [EC1-] Research must be written in English 
\item [EC2-] Research must be related to software engineering or computer science 
\item [EC3-] Research from conference proceedings, book chapters, magazines, news, and posters, or grey literature are excluded.
\item [EC4-] Research must be published from 2013 to 2023. 
\item [EC5-] Research Full-text is available.
\item [EC6-] Research is not duplicated
\end{enumerate}

\subsection{Primary Clustering}
The main advantages and disadvantages of the databases chosen to conduct the search are presented in Table 1. By examining the information presented in this table, it is possible for the reader to determine that all of the databases used are reliable and of high academic quality.

\begin{table}[]
\caption{Table 1 - Databases Pros and Cons}
\begin{tabular}{|l|l|l|}
\hline
Database   & Pros                                                          & Cons                                                                                                                                  \\ \hline
IEEExplore & peer-reviewed                                                       & \begin{tabular}[c]{@{}l@{}}-Require subscription \\ -Some papers can’t be downloaded directly\end{tabular} \\ \hline
ScienceDirec     & \begin{tabular}[c]{@{}l@{}}Filter by\\CS categories\end{tabular} & \begin{tabular}[c]{@{}l@{}}-No advanced search\end{tabular} \\ \hline
ACM        & peer-reviewed                                                       & \begin{tabular}[c]{@{}l@{}}Access to papers only through subscription\\ -Some papers can’t be downloaded directly\end{tabular}                 \\ \hline
\end{tabular}
\end{table}

Figure \ref{cluster_digitalLibraries} displays how the papers are distributed across the selected databases, while Figure \ref{clusterPublishers} shows their distribution across publishers. We found a significant representation of papers from different digital libraries. IEEE Xplore Digital Library stands out as the most utilized platform, hosting 18 out of the total papers reviewed. The Association for Computing Machinery (ACM) Digital Library follows with 9 papers. The ScienceDirect (SD) platform, although less represented, contributed 3 papers to our review. This distribution underscores the dominance of IEEE Xplore and ACM Digital Library as primary sources of research related to AI in education, with ScienceDirect also serving as a notable resource.

We analyzed papers from several different publishers. The majority of these papers, 18 in total, were published by IEEE. The Association for Computing Machinery (ACM) was the second most common publisher, accounting for 9 papers. Other publishers such as Computers \& Education, Computers in Human Behavior, and Procedia Computer Science each had one paper included in our review. These two figures provide evidence that the papers included in our final selection are sourced from reputable computer science conferences or highly ranked journals.

Furthermore, Figure \ref{clusterByYears} exhibits the distribution of papers by year of publication. It's interesting to observe the frequency distribution of the years these papers were published. The most common publication year was 2019, with a total of 8 papers. This was followed by 2021, which accounted for 6 papers. The years 2020 and 2022 each contributed 5 papers to our review. Earlier years such as 2016 and 2018 were represented by 2 papers each. Lastly, the years 2015 and 2017 were the least common, with only one paper each being published in those years. This distribution provides an overview of the recent trends in research on AI applications in education.



\begin{figure}[htbp]
\centerline{\includegraphics[scale=0.5]{figs/cluster_digitalLibraries.png}}
\caption{Databases distribution}
\label{cluster_digitalLibraries}
\end{figure}

\begin{figure}[htbp]
\centerline{\includegraphics[scale=0.5]{figs/clusterPublishers.png}}
\caption{Publishers distribution}
\label{clusterPublishers}
\end{figure}

\begin{figure}[htbp]
\centerline{\includegraphics[scale=0.5]{figs/clusterByYears.png}}
\caption{Years distribution}
\label{clusterByYears}
\end{figure}

% \section{Results and Discussion}
% \label{results}
% \subsection{\textbf{What Problems researchers are trying to solve and what technologies are utilized by researchers to develop Automated Writing Evaluation (AWE) software?}}
% \label{RQ1}

% Our review of the research in the field of Academic Writing in English (AWE) revealed that it encompasses a vast array of problems that researchers aim to solve. To provide readers with a more organized and accessible understanding of these problems, we have categorized them and listed them below. We have also identified instances where multiple studies have addressed the same topic, and have grouped them together for clarity in Table \ref{Table 2} .

% In order to provide readers with a deeper understanding of the problems addressed in our review, we have included a glossary section that defines each problem mentioned below. This section is intended to provide readers with a comprehensive overview of each problem and its context within the field of AWE.

% \begin{itemize}
% \item 
%  Assess the quality of the written text:   \cite{Towards_Context_Aware_Automated}
% \item Authorship detection in text corpora:  \cite{An_Approach_to_Detecting_Text_Autorship}
% \item Collocation and lexical bundle detection: \cite{Extraction_of_Lexical_Bundles}
% \item Collocation error detection: \cite{The_Effects_of_Learner_Errors}
% \item Collocation recommendation system for academic writing: \cite{A_SentenceWide_Collocation_Recommendation_System}
% \item Discourse and Rhetorical moves in text corpora: \cite{Discourse_Classification},\cite{Causal_discourse_analyzer},\cite{Automatic_Identification_of_Discourse},\cite{A_Corporal}
% \item Extraction of the lexico-semantic patterns from text corpora: \cite{The_LexicoSemantic_Pattern_Extraction}
% \item Generate template-like text for literature review: \cite{Text_Generation_with_Diversified}
% \item Grammatical error detection: \cite{Ensemble_MultiChannel_Neural_Networks}, 
% \cite{Automated_Error_Detection_Developing}
% \item Incoherent sentences: \cite{Incoherent_Sentence_Detection_ScientificArticles} 
% \item Predicting vocabulary score based on linguistic features: \cite{Automatic_Student_Writing_Evaluation}
% \item Style transfer between abstract text and conclusion both ways \cite{Unpaired_Abstract_to_Conclusion_TextStyle}
% \end{itemize}


% % Please add the following required packages to your document preamble:
% % \usepackage[table,xcdraw]{xcolor}
% % If you use beamer only pass "xcolor=table" option, i.e. \documentclass[xcolor=table]{beamer}
% \begin{table*}[t]
%   \centering
%   \begin{tabularx}{\textwidth}{|p{1cm}|X|X|}
% \hline
% \textbf{Research}                                               & \textbf{Technologies}                                                                                                                                                                       & \textbf{Problems to solve}                                                        \\ \hline
% \cite{Towards_Context_Aware_Automated}                  & \begin{tabular}[c]{@{}l@{}}Supervision learning \\ (no mention for algoritihm)\end{tabular}                                                                                         & Assess the quality of the written text                         \\ \hline
% \cite{An_Approach_to_Detecting_Text_Autorship}          & \begin{tabular}[c]{@{}l@{}}Minkowski distance, \\ Chevishev distance,\end{tabular}                                                                                                  & authorship detection in text corpses                           \\ \hline
% \cite{Extraction_of_Lexical_Bundles}                    & \begin{tabular}[c]{@{}l@{}}Bernoulli Naive Bayes, \\ Support Vector Machines, \\ Multilayer Perceptron\end{tabular}                                                                 & Collocation and lexical bundle detection                       \\ \hline
% \cite{The_Effects_of_Learner_Errors}                    & POS Tagger, Wordnet                                                                                                                                                                 & Collocation error detection                                    \\ \hline
% \cite{A_SentenceWide_Collocation_Recommendation_System} & \begin{tabular}[c]{@{}l@{}}Stanford core, \\ Dependency parsers, \\ POS, Sentence detector, \\ Stemmer, \\ Wordnet\end{tabular}                                                     & Collocation recommendation system for academic writing         \\ \hline
% \cite{Discourse_Classification}                         & \begin{tabular}[c]{@{}l@{}}n-gram, \\\end{tabular}                                                                                & Discourse and Rhetorical moves in text corpses                 \\ \hline
% \cite{Causal_discourse_analyzer}                        & \begin{tabular}[c]{@{}l@{}}Stanford core, \\ Dependency parsers, \\ POS, Sentence detector, \\ Constitute tree, \\ xml parsing, \\ prolog Programing language, \\ perl\end{tabular} & Discourse and Rhetorical moves in text corpses                 \\ \hline
% \cite{Automatic_Identification_of_Discourse}            & \begin{tabular}[c]{@{}l@{}}n-gram, Move step classifier,\\ Support Vector Machines SVM\end{tabular}                                                                                 & Discourse and Rhetorical moves in text corpses                 \\ \hline
% \cite{A_Corporal}                                       & Antconc software, lda algorithm                                                                                                                                                     & Discourse and Rhetorical moves in text corpses                 \\ \hline
% \cite{The_LexicoSemantic_Pattern_Extraction}            & GATE and OpenNLP                                                                                                                                                                    & Extraction of the lexico-semantic patterns from text corpses   \\ \hline
% \cite{Text_Generation_with_Diversified}                 & Markov model                                                                                                                                                                        & Generate template like text for literature review,             \\ \hline
% \cite{Ensemble_MultiChannel_Neural_Networks}            & EMC-NN, BiLSTM, and CNN                                                                                                                                                             & Grammatical error detection                                    \\ \hline
% \cite{Automated_Error_Detection_Developing}             & \begin{tabular}[c]{@{}l@{}}Stanford core, \\ Dependency parsers, \\ POS, \\ Sentence detector, \\ Constitute tree, \\ xml parsing, \\ prolog Programing language\end{tabular}       & Grammatical error detection                                    \\ \hline
% \cite{Incoherent_Sentence_Detection_ScientificArticles} & \begin{tabular}[c]{@{}l@{}}BERT-base-uncased\\ BERT-large-uncased,\end{tabular}                                                                                                     & Incoherent sentences                                           \\ \hline
% \cite{Automatic_Student_Writing_Evaluation}             & \begin{tabular}[c]{@{}l@{}}Linear Regression, \\ SVM, \\ Linear SVM Polynomial, \\ Random Forest, \\ and Gradient Boosting Tree\end{tabular}                                        & Predicting vocabulary score based on linguistics features      \\ \hline
% \cite{Unpaired_Abstract_to_Conclusion_TextStyle}        & \begin{tabular}[c]{@{}l@{}}cross align, \\ DualRein, \\ Sentence-BERT, \\ Cycle-GAN\end{tabular}                                                                 & Style transfer, between abstract text and conclusion both ways \\ \hline
% \end{tabularx}
% \caption{Problems and technologies used }
% \label{Table 2}
% \end{table*}

% In addition to categorizing the research problems, we also analyzed the technologies utilized by the researchers in their studies. To identify commonalities in the technologies employed, we gathered data from each study and generated a bar chart displaying the frequency of usage. Our analysis reveals that SVM, Dependency parser, and Standford Core were the most frequently used technologies, while POS and BERT were used in a smaller number of studies, but still showed significant usage. The chart provides a clear visualization of the distribution of technology usage among the studies reviewed.

% % \begin{figure}[t]
% %   \centering
% %   \includegraphics[width=\textwidth]{images/technolgies.png}
% %   \caption{Caption of the figure.}
% %   \label{fig:my-figure}
% % \end{figure}

% \begin{figure*}
%     \centering
%     \includegraphics{figs/technolgies.png}
%     \caption{Technologies}
%     \label{fig:arch}
% \end{figure*}



% \subsection{\textbf{RQ2 - What datasets were used in the research?}}
% \label{RQ2}

% In our review of research papers related to AWE, we have encountered a wide range of datasets that have been utilized to address various problems. To provide a comprehensive view, we collected the names of these datasets and any available information, such as the number of articles that have been processed using them. Table 3 lists all of the datasets that we have reviewed. This analysis provides valuable insights for researchers looking for reliable datasets to use in their AWE-related research. Moreover, most of these datasets are publicly available on the internet, making them easily accessible for any interested researchers.


% \begin{table*}[]
%   \centering
%   \begin{tabularx}{\textwidth}{|p{1cm}|X|}
% \hline
% \hline
% \multicolumn{1}{|c|}{\textbf{Research}}                 & \multicolumn{1}{c|}{\textbf{Datasets}}                                                                                                                                                                                                                \\ \hline
% \cite{A_Corporal}                                       & 1,500 abstracts conference papers from Advanced Learning Technologies (lCALT).                                                                                                                                                                        \\ \hline
% \cite{Ensemble_MultiChannel_Neural_Networks}            & AESW datasets, texts extracted from 9,919 selected papers                                                                                 
% \\ \hline
% \cite{The_Effects_of_Learner_Errors}                    & 299 TOEFL essays                                                                                                                                                                                                                                      \\ \hline
% \cite{The_LexicoSemantic_Pattern_Extraction}            & 49 scientific publications on Computer Science domain                                                                                                                                                                                                       \\ \hline
% \cite{Unpaired_Abstract_to_Conclusion_TextStyle}        & ACL-ARC, 19385 abstract and 16317 conclusion                                                                                                                                                                                                          \\ \hline
% \cite{Incoherent_Sentence_Detection_ScientificArticles} & Arxiv,TimeTravel Dataset, and Cyberleninka Dataset                                                                                                                                                                                                    \\ \hline
% \cite{Text_Generation_with_Diversified}                 & \begin{tabular}[c]{@{}l@{}}Dergipark, TRDizin, Council of Higher Education Thesis Center and SOBIAD\\ 100 thousand theses and more than 250 thousand academic publications full text\end{tabular}                                                     \\ \hline
% \cite{An_Approach_to_Detecting_Text_Autorship}          & Database composed of 3000 short articles written in Portuguese.                                                                                                                                                                                       \\ \hline
% \cite{Discourse_Classification}                         & large specialized corpus of 1,020 research articles, and 1,322,089 words                                                                                                                                                                              \\ \hline
% \cite{Extraction_of_Lexical_Bundles}                    & NLP4NLP , ACL Anthology Reference Corpus2 (ACL-ARC hereafter) with 21,636 articles.                                                                                                                                                                   \\ \hline
% \cite{Towards_Context_Aware_Automated}                  & No mention                                                                                                                                                                                                                                            \\ \hline
% \cite{Automated_Error_Detection_Developing}             & \begin{tabular}[c]{@{}l@{}}Privtate data, (papers submitted from nonnative students in English department) not all of the papers are scintific\\ 5500 sentece randomly sampled togther with corrective feedback from criterion\end{tabular}           \\ \hline
% \cite{Automatic_Identification_of_Discourse}            & The corpus used in this study was compiled from an existing corpus of published research articles in 44 disciplines, used in an academic writing graduate course for international students. The corpus con- tains 1,623 articles and 1,322,089 words \\ \hline
% \cite{Causal_discourse_analyzer}                        & 585 sentences total in 17 essays in a varional topics                                                                                                                                                                                                 \\ \hline
% \cite{Automatic_Student_Writing_Evaluation}             & 106 writing eassays                                                                                                                                                                                                                                   \\ \hline
% \cite{A_SentenceWide_Collocation_Recommendation_System} & 3000 SCI articles in journals of engineering and science fields                                                                                                                                                                                       \\ \hline
% \end{tabularx}
% \caption{Datasets used in the researches}
% \label{Table 3}
% \end{table*}



% \subsection{\textbf{What methods were employed to validate the accuracy of the AWE software models and what are the values for this evaluation?}}
% \label{RQ3}


% During our research, we thoroughly analyzed the papers to determine how each proposed solution was evaluated. Most of the papers we reviewed documented the evaluations performed during the development of their solutions or models to solve the problems at hand. We gathered all of the evaluation values and the metrics used in Table \ref{Table 4}. This table is particularly beneficial to researchers who wish to quickly view the achieved performance and the problem addressed for each evaluation. We made sure to include all relevant information to ensure maximum usefulness for the reader.


% \begin{table*}[]
% \begin{tabularx}{\textwidth}{|p{1cm}|X|X|}
% \hline
% \multicolumn{1}{|c|}{\textbf{Research}}                                                                                                                                                       & \multicolumn{1}{c|}{\textbf{Evaluation and Metrics}}                                                                                                                                                                   & \multicolumn{1}{c|}{\textbf{Problem research are solving}}                                                                                                                                                                                                            \\ \hline
% \cite{A_Corporal}, \cite{The_Effects_of_Learner_Errors}, \cite{The_LexicoSemantic_Pattern_Extraction} ,\cite{An_Approach_to_Detecting_Text_Autorship}, \cite{Towards_Context_Aware_Automated} & No mention for evaluation                                                                                                                                                                                                             & \begin{tabular}[c]{@{}l@{}}- Discourse and Rhetorical moves in text corpses, \\ - Collocation error detection\\ - Extraction of the lexico-semantic patterns from text corpses\\ - authorship detection in text corpses\\ - Assess the quality of the written text\end{tabular} \\ \hline
% \cite{Automatic_Student_Writing_Evaluation}                                                                                                                                                   & \begin{tabular}[c]{@{}l@{}}RMSE for the following models:\\ Linear Regression = 19.199, \\ Polynomial SVM= 17.302, \\ Linear SVM = 18.502, \\ Random Forest 16.650, \\ and Gradient Boosting Tree= 17.117\end{tabular} & Predicting vocabulary score based on linguistics features                                                                                                                                                                                                             \\ \hline
% \cite{Extraction_of_Lexical_Bundles}                                                                                                                                                          & \begin{tabular}[c]{@{}l@{}}Accuracy = 0.762, \\ Precision = 0.751, \\ Recall = 0.701, \\ F-Score = 0.725\end{tabular}                                                                                                  & Collocation and lexical bundle detection                                                                                                                                                                                                                              \\ \hline
% \cite{Unpaired_Abstract_to_Conclusion_TextStyle}                                                                                                                                              & \begin{tabular}[c]{@{}l@{}}Accuracy = 75.84\\ BLUE=27.23\\ Perplexity=19.57\end{tabular}                                                                                                                               & Style transfer, between abstract text and conclusion both ways                                                                                                                                                                                                        \\ \hline
% \cite{Automatic_Identification_of_Discourse}                                                                                                                                                  & Interannotator agreement K = 0.984                                                                                                                                                                                     & Discourse and Rhetorical moves in text corpses                                                                                                                                                                                                                        \\ \hline
% \cite{Ensemble_MultiChannel_Neural_Networks}                                                                                                                                                  & \begin{tabular}[c]{@{}l@{}}Precision = 0.5378, \\ Recall= 0.7803, \\ F1 = 0.6367\end{tabular}                                                                                                                          & Grammatical error detection                                                                                                                                                                                                                                           \\ \hline
% \cite{Discourse_Classification}                                                                                                                                                               & \begin{tabular}[c]{@{}l@{}}Precision: 70\%, Recall, 61\%, F1 65\% - Move classifier \\ Precision: 68\%, Recall, 54.9\%, F1 61\% - Step classifier\end{tabular}                                                         & Discourse and Rhetorical moves in text corpses                                                                                                                                                                                                                        \\ \hline
% \cite{Incoherent_Sentence_Detection_ScientificArticles}                                                                                                                                       & \begin{tabular}[c]{@{}l@{}}Precision= 0.594, \\ Recall =0.431, \\ F1 = 0.500, \\ and Accuracy 0.712\end{tabular}                                                                                                       & Incoherent sentences                                                                                                                                                                                                                                                  \\ \hline
% \cite{Text_Generation_with_Diversified}                                                                                                                                                       & \begin{tabular}[c]{@{}l@{}}Presceion: 0,78\\ Recall 0.893\\ Accuracy 0.800\\ F-mesure 0.833\end{tabular}                                                                                                               & Generate template like text for literature review,                                                                                                                                                                                                                    \\ \hline
% \cite{Causal_discourse_analyzer}                                                                                                                                                              & \begin{tabular}[c]{@{}l@{}}Precision = 0.83, \\ Recall = 0.84, \\ F1 = 0.83, \\ Accuracy = 0.90\end{tabular}                                                                                                           & Discourse and Rhetorical moves in text corpses                                                                                                                                                                                                                        \\ \hline
% \cite{Automated_Error_Detection_Developing}                                                                                                                                                   & \begin{tabular}[c]{@{}l@{}}CyWrite articles evaluation: \\ P = 0.56\\ R = 0.76\\ F= 0.65\\ Accuracy = 0.84\end{tabular}                                                                                                & Grammatical error detection                                                                                                                                                                                                                                           \\ \hline
% \cite{A_SentenceWide_Collocation_Recommendation_System}                                                                                                                                       & Manually by Expertes in the English domin 73\%                                                                                                                                                                         & Collocation recommendation system for academic writing                                                                                                                                                                                                                \\ \hline
% \end{tabularx}
% \caption{Evaluation and Metrics}
% \label{Table 4}
% \end{table*}

% Worth to mention in Table \ref{Table 5}, we included all research that has created an AWE software, and talked about its access states.

% \begin{table}[]
% \begin{tabular}{|l|l|l|}
% \hline
% \multicolumn{1}{|c|}{\textbf{Research}}                 & \multicolumn{1}{c|}{\textbf{Name}} & \multicolumn{1}{c|}{\textbf{Status}} \\ \hline
% \cite{Automatic_Student_Writing_Evaluation}             & AWT                                & under development                    \\ \hline
% \cite{Automated_Error_Detection_Developing}             & CyWrite                            & Available as open source on Github   \\ \hline
% \cite{A_SentenceWide_Collocation_Recommendation_System} & No mention             & Not accesable                        \\ \hline
% \end{tabular}
% \caption{Tools developed by Researches}
% \label{Table 5}
% \end{table}







	 


\section{Results and Discussion}
\label{Results and Discussion}

\subsection{\textbf{RQ1.  What are the use cases the AI is used to gain insight to how enhance the education? and what are the AI algorithms/techniques used for such purpose?}} \label{RQ1_analysis}


\begin{flushleft}
After conducting a thorough review of the papers, we have identified five distinct use cases for the research:
\end{flushleft}

\begin{enumerate}
\item \textbf{Use Case 1-Trio-Techniques Text Analysis for Focused Educational Feedback:} papers which used Sentiment Analysis, Topic Modeling, and Text Clustering Combined
\item \textbf{Use Case 2 - Duo-Technique Text Analysis for Focused Educational Feedback:} papers which used Sentiment Analysis and Topic Modeling Together or Sentiment Analysis and Clustering Together
\item \textbf{Use Case 3- Solo-Technique Text Analysis for Specific Educational Enhancements: }papers which used Sentiment Analysis alone, Text Clustering alone, or Topic Modeling alone
\item \textbf{Use Case 4- Neural Network-Based Forecasting for Educational Outcomes:} Neural Network for Quantitative Data Prediction
\item \textbf{Use Case 5 - Computer Vision for tracking emotions:} paper which used Computer Vision approach. 
\end{enumerate}

Among these use cases, \cite{Palaute:_An_Online_Text_Mining_Tool_for_Analyzing_Written_Student_Course_Feedback} falls under \textbf{Use Case 1} they used Sentiment Analysis  to see sentiment values related to course ratings, difficulty levels, and learning duration. Moreover, Emotion Analysis is executed on student feedback, providing an intriguing layer of emotional context to the raw data. The paper also utilizes Topic Modeling to categorize and analyze the key themes present in student feedback. The findings are effectively visualized using color-coded circles to denote sentiment, aiding in data interpretation. This combination of AI methodologies offers a comprehensive approach to understanding and enhancing educational experiences based on students feedback.

\textbf{Use Case 2}, \cite{Analyzing_educational_comments_for_topics_and_sentiments:_A_text_analytics_approach} used text Classification to distinguish between a variety of topics including (Content, Technology, Interaction, Pace, Experience), then calculated sentiment analysis for each of the classified topics. Lastly, they did clustering to see which are the hot topics that the students talk about the most, and in which sentiment—positive or negative. In \cite{Visualizing_Student_Opinion_Through_Text_Analysis} The researchers in the paper employed a systematic approach to analyze satisfaction versus comments in their dataset. First, the data processing, which included converting to lower case, removing stop words, punctuation, The researchers then used Latent Dirichlet Allocation (LDA) for topic modeling.
Comments containing common conjunctions were split into separate comments to avoid the confusion that can arise when a single sentence carries both positive and negative sentiments. For instance, a comment might state, "the pictures are good but the tutorials are bad".
AFINN was employed for sentiment analysis, providing a range between -5 and +5 for more precise sentiment detection. Finally, the researchers utilized visualizations to display sentiment related to specific topic or average satisfaction. The latter was calculated by summing the satisfaction of all comments relating to a particular aspect, then dividing by the total number of such comments. This meticulous approach ensures more accurate and nuanced data interpretation.


Falls under \textbf{Use Case 3}, utilizing Topic modeling alone the paper \cite{A_Study_of_MOOC_Course_Review_Topics_Mining_Based_on_LDA_Topic_Model} presents a comprehensive methodology for analyzing MOOC data on computer courses. Initially, data was collected through web crawling of MOOC platforms. The subsequent data processing phase involved steps such as data de-duplication, Jieba participle, and removal of deactivated words like "very good" and "but".

They utilized Latent Dirichlet Allocation (LDA) to extract topics of user interest, with eight topics being selected from the LDA topic modeling. Visualization of these results was done using pyLDAvis, enhancing interpretability of the findings.

Additionally, a word cloud was generated featuring the top 100 words with the highest frequency of occurrence, providing an immediate visual summary of the most discussed subjects. This combination of techniques serves to effectively explore and display key areas of user attention within the computer courses.

falls under \textbf{Use Case 4}, paper \cite{BP_Neural_Network-Based_Evaluation_on_University_Teachers'_Teaching_Quality} aims to create a robust, scientific system to evaluate teachers' quality using multiple evaluation indices such as teaching attitude, course content, and effectiveness, among others. These indices are inputs for secondary systems, whose outputs feed into a neural network, ultimately determining the final teaching quality evaluation. 

Each secondary system employs a three-layer neural network, with the weights and thresholds determined through sample training. The comprehensive evaluation system uses a similar three-layer Backpropagation (BP) neural network, taking outputs from each secondary system as inputs. 

The final output is divided into five grades (excellent, good, medium, pass, fail) corresponding to specific score ranges. Thus, the system provides a detailed, layered, and reliable evaluation of teaching quality.

Lastly, falls under \textbf{Use Case 5} 
The paper focuses on analyzing the relationship between learners' facial expressions and teaching effectiveness in an online environment. It establishes a real-time facial expression recognition system that tracks and analyzes the learners' state, assisting teachers in evaluating the effectiveness of online teaching and controlling the teaching process.


Through observing the changes in learners' facial expressions, the paper aims to understand their emotional state and thereby evaluate the impact of online learning. This involves processes like face detection, extraction of facial expression features, and the application of facial expression classification algorithms. The study utilizes Convolutional Neural Networks (CNN) in TensorFlow to achieve these objectives.

It is worth mentioning the work of Palaute et al. \cite{Palaute:_An_Online_Text_Mining_Tool_for_Analyzing_Written_Student_Course_Feedback}, which investigates sentiment and emotional analysis to gain insights into students' mental health. Another relevant study by SUFAT \cite{SUFAT_An_Analytics_Tool_for_Gaining_Insights_from_Student_Feedback_Comments} focuses on extracting suggestions from student feedback comments to enhance the overall quality of education.

In conclusion, all of these use cases have a direct impact on enhancing the quality of education through data analysis, whether it involves quantitative or qualitative data. Table \ref{Use Cases and Paper References} in our paper provides a distribution of these use cases across the reviewed papers. Please note that this table assumes the use cases are mutually exclusive, meaning that a paper fitting into Use Case 1 cannot be listed under Use Case 2 or 3, and so on.


\begin{table}[htbp]
\centering
\caption{Use Cases and Paper References - QDP: Quantitative data prediction, SA: Sentiment analysis, TM: Topic modeling, TC: Text clustering, CV: Computer vision, UC: Use case }
\label{Use Cases and Paper References}
\begin{tabular}{|l|p{0.45\linewidth}|}
\hline
\textbf{Use Case} & \textbf{Paper References} \\ \hline
UC-1: SA+TM+TC & \cite{An_Efficient_Data_Mining_Technique_for_Assessing_Satisfaction_Level_With_Online_Learning_for_Higher_Education_Students_During_the_COVID-19} \\ \hline
UC-2: SA + TA OR SA + TC & \cite{Palaute:_An_Online_Text_Mining_Tool_for_Analyzing_Written_Student_Course_Feedback}, \cite{Visualizing_Student_Opinion_Through_Text_Analysis}, \cite{Experience_in_learners_review_to_determine_attribute_relation_for_course_completion},
\cite{SUFAT_An_Analytics_Tool_for_Gaining_Insights_from_Student_Feedback_Comments}, \cite{Understanding_Students_Online_Reviews_to_Improve_College_Experience_and_Graduation_Rates_of_STEM_Programs_at_the_Largest_Post-Secondary_Institution_A_Learner-Centered_Study}, \cite{Assessing_learners_satisfaction_in_collaborative_online_courses_through_a_big_data_approach},\cite{Analyzing_educational_comments_for_topics_and_sentiments:_A_text_analytics_approach} \\ \hline
UC-3: SA OR TA OR TC & \cite{Sentiment_analysis_of_student_feedback_using_machine_learning_and_lexicon-based_approaches}, \cite{A_Study_of_MOOC_Course_Review_Topics_Mining_Based_on_LDA_Topic_Model}, \cite{What_predicts_student_satisfaction_with_MOOCs:_A_gradient_boosting_trees_supervised_machine_learning_and_sentiment_analysis_approach}, \cite{Research_on_Project-Based_Teaching_Quality_Evaluation_Based_on_Fuzzy_Cluster_Analysis}, \cite{Improving_MOOC_quality_using_learning_analytics_and_tools}, \cite{Course_Satisfaction_in_Engineering_Education_Through_the_Lens_of_Student_Agency_Analytics}, \cite{Personalized_Learning_Model_based_on_Deep_Learning_Algorithm_for_Student_Behaviour_Analytic}, \cite{Online_reviews_evaluation_system_for_higher_education_institution:_An_aspect_based_sentiment_analysis_tool} \\ \hline
UC-4: QDP & \cite{Competition_Neural_Network_Model_Application_in_the_Bilingual_Teaching_Quality_Evaluation_of_Universities}, 
\cite{Romanov}\cite{Research_on_the_Application_of_BP_Neural_Network_in_Vocal_Music_Teaching_Quality_Evaluation}, \cite{Design_of_online_teaching_quality_evaluation_system_for_Private_University_Research_based_on_deep_learning_algorithm}, \cite{BP_Neural_Network-Based_Evaluation_on_University_Teachers'_Teaching_Quality}, \cite{An_Improved_BP_Neural_Network-Based_Quality_Evaluation_Model_for_Chinese_International_Education_Teaching_Courses}, \cite{Evaluation_Model_of_PE_Teaching_Quality_Based_on_Active_Learning_Support_Vector_Machine}, \cite{Evaluation_System_of_Tourism_Psychology_Teaching_Quality_Based_on_Convolutional_Neural_Network}, \cite{An_Evaluation_Model_of_Moral_Education_Teaching_Quality_in_Universities_Based_on_System_Control_Theory}, \cite{Research_on_the_Application_of_BP_Neural_Network_in_the_Evaluation_of_College_Teaching_Quality}, \cite{The_Evaluation_and_Application_of_BP_Neural_Network_in_the_Multiple_Quality_Assurance_Paths_of_Higher_Education}, \cite{Evaluating_the_quality_of_teaching_in_higher_education_institutes_using_clustering_approach_and_set_pair_analysis}, \cite{Automatic_Rules_Generation_for_Teaching_Strategies_Improvement}, \cite{Exploring_Causes_for_the_Dropout_on_Massive_Open_Online_Courses} \\ \hline
UC-5: CV & \cite{Research_on_Evaluation_of_Online_Teaching_Effect_based_on_Deep_Learning_Technology} \\ \hline
\end{tabular}
\end{table}


The analysis of the findings reveals a diverse range of AI techniques and applications employed in measuring and enhancing the quality of education illustrated in figure \ref{Ai_teqniuqes}. Among these, sentiment analysis emerges as the most frequently utilized technique, appearing 26 times in the findings. By analyzing and interpreting the sentiment expressed in educational data, researchers gain valuable insights into students' attitudes, perceptions, and emotions.
Text classification, text clustering, and topic modeling are also prominent techniques, with varying frequencies of occurrence. These techniques enable researchers to organize, group, and uncover latent themes within educational texts, providing a deeper understanding of the content and underlying concepts.

Quantitative data prediction appearing nine times in the findings. This technique allows researchers to utilize quantitative data to predict various outcomes in education, such as dropout rates, student performance metrics, and satisfaction scores. By leveraging predictive models, educators and policymakers can gain insights into trends and patterns that impact the quality of education

Feature selection, emotional analysis, data processing, classification, also demonstrate their relevance in the field. These techniques contribute to refining data, analyzing emotions, making data-driven decisions.

The findings underscore the diverse nature of AI techniques applied in education and highlight their potential in advancing the quality of education through data analysis, prediction, and classification. They provide researchers and practitioners valuable insights into the prevailing practices and avenues for further exploration in the quest for improving educational outcomes.

\begin{figure}[htbp]
\centerline{\includegraphics[scale=0.6]{figs/Ai_teqniuqes.png}}
\caption{AI techniques used across papers}
\label{Ai_teqniuqes}
\end{figure}

Figure \ref{Ai_algorthim} represents the frequency of various AI techniques used in the field of education. Among the techniques, 'SVM' (Support Vector Machine) is the most commonly employed, appearing 6 times. 'Scikit-learn', 'BP NN' (Backpropagation Neural Network), and 'MATLAB' follow closely with frequencies of 4, 6, and 3 respectively. Other techniques such as 'NLTK', 'Spacy', 'NRC', 'Syuzhet', 'CNN', 'TF-IDF', 'n-gram', 'Naïve Bayes', 'AFINN', 'Random Forest', 'K-Means', and 'LDA' also find notable usage. These AI techniques showcase the diversity of approaches employed in the analysis and processing of educational data, enabling researchers to extract meaningful insights and improve the quality of education through advanced machine learning algorithms and natural language processing.

Lastly, we noticed a varying distribution across different types of educational settings ilustrated in figure \ref{clusterEducationType}. A substantial majority of the papers, 15 in total, pertained to the field of Online Higher Education, indicating a strong focus on the application of AI in virtual tertiary level institutions. Traditional Higher Education was the second most discussed setting, with 9 papers reflecting its relevance. Lastly, MOOCs (Massive Open Online Courses), an emerging platform for online learning, was the subject of 5 papers. This distribution underscores the increasing interest in leveraging AI for improving online education and the broader higher education sector

\begin{figure}[htbp]
\centerline{\includegraphics[scale=0.50]{figs/clusterEducationType.png}}
\caption{Education type distribution}
\label{clusterEducationType}
\end{figure}

\begin{figure}[htbp]
\label{Ai_algorthim}
\centerline{\includegraphics[scale=0.50]{figs/Ai_algorthims.png}}
\caption{AI algorithms used across papers}
\label{fig}
\end{figure}




\subsection{\textbf{RQ2. What are the metrics commonly used to assess the quality of online and offline education, what are the performance factors they used to evaluate theses metrics?}} \label{RQ2_analysis}

Out of the papers we reviewed we noticed that researchers frequently utilize various terms and phrases to express their methods of measuring education quality as metrics. Despite differences in wording, many of these terms are semantically similar. For instance, two distinct studies might use the metrics 'Teaching quality' and 'Teaching quality evaluation'; though expressed differently, their meanings are essentially the same. We have therefore categorized these various metrics into four broad groups for clarity and conciseness. These categories are: Teaching Quality and Effectiveness, Teaching Evaluation and Feedback, Learner Satisfaction and Experience, and Course Completion and Retention. The specific metrics that fall under each category are illustrated in Table \ref{Table metrics categories}.



\begin{table}[h!]
\centering
\caption{Categories and Their Corresponding Metrics}
\label{Table metrics categories}
\begin{tabular}{|p{0.15\textwidth}|p{0.35\textwidth}|}
\hline
\textbf{Category} & \textbf{Metrics} \\ \hline
Teaching Quality and Effectiveness & Teaching quality, Teacher's teaching quality, Satisfaction with the quality of teaching, Teaching quality of the curriculum, Project-based teaching quality, MOOCs quality, Teaching quality evaluation \\ \hline
Teaching Evaluation and Feedback & Student evaluations of teaching (SET), Teaching evaluating system, Teaching evaluation system TES, Student performance \\ \hline
Learner Satisfaction and Experience & Student satisfaction level, Students satisfaction, Learning experience, MOOC learner satisfaction, Learner satisfaction metric, Students satisfaction score \\ \hline
Course Completion and Retention & Learners' course completion, MOOCs dropout \\ \hline
\end{tabular}
\end{table}


\begin{table*}
\centering
\caption{Categories and Their Corresponding Metrics}
\label{Categories metrics}
\begin{tabular}{|p{0.22\textwidth}|p{0.3\textwidth}|p{0.22\textwidth}|}
\hline
\textbf{Category} & \textbf{Metrics} & \textbf{Paper References} \\ \hline
Teaching Quality and Effectiveness & Teaching quality, Teacher's teaching quality, Satisfaction with the quality of teaching, Teaching quality of the curriculum, Project-based teaching quality, MOOCs quality, Teaching quality evaluation & 
\cite{Competition_Neural_Network_Model_Application_in_the_Bilingual_Teaching_Quality_Evaluation_of_Universities},
\cite{Visualizing_Student_Opinion_Through_Text_Analysis}\cite{Improving_MOOC_quality_using_learning_analytics_and_tools},
\cite{Research_on_the_Application_of_BP_Neural_Network_in_Vocal_Music_Teaching_Quality_Evaluation}, \cite{Evaluating_the_quality_of_teaching_in_higher_education_institutes_using_clustering_approach_and_set_pair_analysis}, \cite{An_Improved_BP_Neural_Network-Based_Quality_Evaluation_Model_for_Chinese_International_Education_Teaching_Courses}, \cite{Research_on_Project-Based_Teaching_Quality_Evaluation_Based_on_Fuzzy_Cluster_Analysis}, \cite{Improving_MOOC_quality_using_learning_analytics_and_tools}, \cite{Online_reviews_evaluation_system_for_higher_education_institution:_An_aspect_based_sentiment_analysis_tool}, \cite{Automatic_Rules_Generation_for_Teaching_Strategies_Improvement}, \cite{Analyzing_educational_comments_for_topics_and_sentiments:_A_text_analytics_approach}, \cite{An_Evaluation_Model_of_Moral_Education_Teaching_Quality_in_Universities_Based_on_System_Control_Theory}, \cite{Research_on_the_Application_of_BP_Neural_Network_in_the_Evaluation_of_College_Teaching_Quality}, 
\cite{The_Evaluation_and_Application_of_BP_Neural_Network_in_the_Multiple_Quality_Assurance_Paths_of_Higher_Education},
\cite{Research_on_Evaluation_of_Online_Teaching_Effect_based_on_Deep_Learning_Technology},
\cite{The_Evaluation_and_Application_of_BP_Neural_Network_in_the_Multiple_Quality_Assurance_Paths_of_Higher_Education} \\ \hline
Teaching Evaluation and Feedback & Student evaluations of teaching (SET), Teaching evaluating system, Teaching evaluation system TES, Student performance & 
\cite{Palaute:_An_Online_Text_Mining_Tool_for_Analyzing_Written_Student_Course_Feedback}
, \cite{Romanov}
\cite{Evaluation_System_of_Tourism_Psychology_Teaching_Quality_Based_on_Convolutional_Neural_Network},
\cite{Personalized_Learning_Model_based_on_Deep_Learning_Algorithm_for_Student_Behaviour_Analytic} \\ \hline
Learner Satisfaction and Experience & Student satisfaction level, Students satisfaction, Learning experience, MOOC learner satisfaction, Learner satisfaction metric, Students satisfaction score, Course satisfaction & 
\cite{Exploring_Causes_for_the_Dropout_on_Massive_Open_Online_Courses},
\cite{An_Efficient_Data_Mining_Technique_for_Assessing_Satisfaction_Level_With_Online_Learning_for_Higher_Education_Students_During_the_COVID-19},
\cite{Visualizing_Student_Opinion_Through_Text_Analysis},
\cite{What_predicts_student_satisfaction_with_MOOCs:_A_gradient_boosting_trees_supervised_machine_learning_and_sentiment_analysis_approach},
\cite{Course_Satisfaction_in_Engineering_Education_Through_the_Lens_of_Student_Agency_Analytics},
\cite{Assessing_learners_satisfaction_in_collaborative_online_courses_through_a_big_data_approach},
\cite{Understanding_Students_Online_Reviews_to_Improve_College_Experience_and_Graduation_Rates_of_STEM_Programs_at_the_Largest_Post-Secondary_Institution_A_Learner-Centered_Study} \\ \hline
Course Completion and Retention & Learners' course completion, MOOCs dropout & 
\cite{Exploring_Causes_for_the_Dropout_on_Massive_Open_Online_Courses},
\cite{Experience_in_learners_review_to_determine_attribute_relation_for_course_completion} \\ \hline
\end{tabular}
\end{table*}




\label{metricsCategories}
\begin{figure}[htbp]
\centerline{\includegraphics[scale=0.50]{figs/metricsCategories.png}}
\caption{Metrics Categories.}
\label{metricsCategories}
\end{figure}


In our analysis, we compared the usage of various metrics to discern research trends in AI-enhanced education quality. This exploration helps identify key focus areas and the importance attributed to each metric. It provides insights into prevailing academic interests, potential gaps in the research, and guides future studies for a balanced and comprehensive investigation of the field. The comparison is better represented in figure \ref{metricsCategories}. below is the interpretation about the metrics usage we analysed and the complete distrbution of the paper can be found in Table \ref{Categories metrics}:
















\begin{enumerate}
\item \textbf{Teaching Quality and Effectiveness:} This category appears to have the highest metric usage with a count of 7. It seems to be the most popular category, which may suggest that most of the focus in the data is on evaluating the effectiveness of the teaching quality and how to improve it.
\item \textbf{Teaching Evaluation and Feedback:} The metric usage for this category is moderately high with a count of 4. This indicates that the evaluation and feedback of teaching is also a critical aspect, though not as prominent as teaching quality and effectiveness.
\item \textbf{Learner Satisfaction and Experience:} This category has the second highest usage with a count of 6. It suggests that a significant emphasis is also given to the satisfaction and experience of the learners. This might be due to the recognition of its impact on overall learning outcomes.
\item \textbf{Course Completion and Retention:} This category has the lowest usage with a count of 2, indicating that less focus is put on course completion and retention compared to the other categories. This might be because other aspects of teaching and learning, such as teaching quality and learner satisfaction, are perceived to be more influential on overall education quality.
\end{enumerate}

lastly, we should mention that papers\cite{Evaluation_Model_of_PE_Teaching_Quality_Based_on_Active_Learning_Support_Vector_Machine} and \cite{Sentiment_analysis_of_student_feedback_using_machine_learning_and_lexicon-based_approaches} didn't mention any metrics.

Based on the usage counts, most of the focus is on `Teaching Quality and Effectiveness' followed by `Learner Satisfaction and Experience'. `Teaching Evaluation and Feedback' and `Course Completion and Retention' also play roles but are less emphasized in the current data. It could be beneficial to consider whether a more balanced focus across all these categories might lead to a more comprehensive understanding of the quality of education.

The second phase of our metric analysis involves examining the factors utilized in the computation of these metrics. Based on the findings, we can structure and categorize the various factors of educational quality into four main groups listed below. Each of them mention categories encompasses several factors, all of which are detailed in Table \ref{Metrics factors}:

\begin{enumerate}
    \item \textbf{Educator Performance Metrics:} This group focuses on factors related to teaching quality, style, and effectiveness.
    \item \textbf{Student Engagement and Performance Metrics:} This group measure student behaviors and achievement, which also indirectly reflect the effectiveness of teaching.
    \item \textbf{Course and Material Quality Metrics:} This group involves the structural components and delivery of the course itself.
    \item \textbf{Environmental and Institutional Metrics:} This group is based on factors that are less directly related to the teacher-student interaction, but still have an influence on educational quality.
\end{enumerate}



\begin{table}[h]
\caption{Metrics Factors}
\label{Metrics factors}
\begin{tabular}{p{0.95\columnwidth}}
\hline
\textbf{Educator Performance Metrics}: Teaching attitude, Teaching method, Teaching content, Teaching ability, Teaching style and ethics, Faculty interaction, Faculty feedback, Faculty preparation, Classroom delivery \\
\hline
\textbf{Student Engagement and Performance Metrics}: Student satisfaction, Student attendance, Student marks, Student achievements, Age, Course rating, Would take again, Login frequency, intervals, and durations \\
\hline
\textbf{Course and Material Quality Metrics}: Course difficulty, Course content, Course description, Courseware design, Learning management, Interactive program, Course structure and contents, Technology and Support, E-learning pace, Overall experience, Course schedule, Course major, Course duration, Time investment required \\
\hline
\textbf{Environmental and Institutional Metrics}: Class style, Teaching staff, School conditions, Quality of students, Campus culture, School spirit, Geographic and psychological environment \\
\hline
\end{tabular}
\end{table}



In conclusion, we analyzed the metrics commonly used to assess the quality of online and offline education. We categorized these metrics, compared their usage, and explored the performance factors associated with them. We also grouped the factors into three main categories to provide a structured overview. Our aim was to gain insights into the evaluation methods used in assessing educational quality.




\subsection{\textbf{RQ-3 What is the impact of AI on the education quality for both online and offline teaching formats?}}
 \label{RQ3_analysis}

In our review of the selected papers, we examined the educational elements impacted by AI techniques and noticed a degree of similarity between these elements. Consequently, we have categorized them into the following clusters:

\begin{itemize}
\item \textbf{Cluster-1: Feedback Analysis and Course Design}
This cluster includes papers which focus on analyzing student feedback and modifying course design accordingly.

\item \textbf{Cluster-2: Teaching Quality and Effectiveness}
Papers in this cluster are centered around evaluating teaching quality and effectiveness, including course content, teacher's effectiveness, and teaching strategies.

\item \textbf{Cluster-3: Course Assessment and Student Satisfaction}
This cluster focuses on course assessment, student satisfaction, and the factors influencing them.

\item \textbf{Cluster-4: Course Duration, Difficulty and Completion Rates}
This cluster includes papers that focus on course duration, difficulty, completion rates, and dropout rates.
\end{itemize}

In the "Cluster-1: Feedback Analysis and Course Design", Palaute et al. \cite{Palaute:_An_Online_Text_Mining_Tool_for_Analyzing_Written_Student_Course_Feedback} introduced a tool designed to extract and analyze student feedback. This tool identified specific areas requiring improvement within the course, thus profoundly influencing the design and faculty-related aspects of the course. Similarly, SUFAT, an analytics tool employed by \cite{SUFAT_An_Analytics_Tool_for_Gaining_Insights_from_Student_Feedback_Comments}, analyzed qualitative feedback comments from students. This information directed instructors toward areas of concern, guiding informed decisions for effective course modifications, thereby enhancing teaching effectiveness and decision-making processes.

In the "Cluster-2: Teaching Quality and Effectiveness", the study by \cite{Research_on_Project-Based_Teaching_Quality_Evaluation_Based_on_Fuzzy_Cluster_Analysis} elucidated students' perceptions of teaching attitude, course content, methods, and effectiveness, consequently identifying areas for improvement. This insight considerably impacted teaching effectiveness and course content. Moreover, the unique approach adopted by \cite{Evaluating_the_quality_of_teaching_in_higher_education_institutes_using_clustering_approach_and_set_pair_analysis} fused student characteristics and feedback to evaluate teaching quality in higher education institutions, thereby influencing the quality of teaching and course evaluations directly.

Moving on to the "Cluster-3: Course Assessment and Student Satisfaction", \cite{An_Efficient_Data_Mining_Technique_for_Assessing_Satisfaction_Level_With_Online_Learning_for_Higher_Education_Students_During_the_COVID-19} proposed a data mining technique to predict and enhance student satisfaction, allowing early interventions to mitigate educational issues, especially during crises such as COVID-19. This intervention directly improved student satisfaction and the efficacy of online learning. Furthermore, \cite{What_predicts_student_satisfaction_with_MOOCs:_A_gradient_boosting_trees_supervised_machine_learning_and_sentiment_analysis_approach} found that aspects such as course assessment, instructor presentation, and course content significantly influenced learner satisfaction, whereas course video and interaction did not significantly impact it.

Lastly, in the "Cluster-4: Course Duration, Difficulty, and Completion Rates" cluster, the research conducted by \cite{Competition_Neural_Network_Model_Application_in_the_Bilingual_Teaching_Quality_Evaluation_of_Universities} identified curriculum quality, class hours, and curriculum difficulty as the most significant factors affecting classification, directly impacting curriculum quality, course duration, difficulty, and completion rates. Simultaneously, \cite{Experience_in_learners_review_to_determine_attribute_relation_for_course_completion} applied classification models to course attributes, revealing correlations between course completion rate, duration, and difficulty.

In conclusion, the research papers provided key insights into how AI techniques can be utilized to improve various elements of education. The effectiveness of AI tools in enhancing course design, teaching quality, student satisfaction, and understanding of course duration, difficulty, and completion rates was evident. For a comprehensive understanding and in-depth analysis of the individual contributions of each research paper, please refer to Table \ref{tab:AI_education_evaluation}.

\begin{table*}
    \centering
    \caption{AI-Enhanced Education Evaluation}
    \label{tab:AI_education_evaluation}
    \begin{tabularx}{\textwidth}{|X|X|X|X|}
        \hline
        \textbf{Cluster} & \textbf{Improved Educational Element} & \textbf{Paper References} \\
        \hline
        Feedback Analysis and Course Design & Course Design \& Faculty, Course modifications, teaching effectiveness, instructor's decision making, Faculty interaction, course structure, Course development, student feedback. & \cite{Palaute:_An_Online_Text_Mining_Tool_for_Analyzing_Written_Student_Course_Feedback}, \cite{SUFAT_An_Analytics_Tool_for_Gaining_Insights_from_Student_Feedback_Comments}, \cite{Analyzing_educational_comments_for_topics_and_sentiments:_A_text_analytics_approach}, \cite{Improving_MOOC_quality_using_learning_analytics_and_tools}, \cite{Visualizing_Student_Opinion_Through_Text_Analysis}, \cite{Online_reviews_evaluation_system_for_higher_education_institution:_An_aspect_based_sentiment_analysis_tool} \\
        \hline
        Teaching Quality and Effectiveness & Curriculum quality, course duration and difficulty, completion rates, Teaching quality, course evaluations, Teaching effectiveness, course content, College teaching quality, Chinese International Education quality. & \cite{Competition_Neural_Network_Model_Application_in_the_Bilingual_Teaching_Quality_Evaluation_of_Universities}, \cite{Evaluating_the_quality_of_teaching_in_higher_education_institutes_using_clustering_approach_and_set_pair_analysis}, \cite{Research_on_Project-Based_Teaching_Quality_Evaluation_Based_on_Fuzzy_Cluster_Analysis}, \cite{Research_on_the_Application_of_BP_Neural_Network_in_the_Evaluation_of_College_Teaching_Quality}, \cite{An_Improved_BP_Neural_Network-Based_Quality_Evaluation_Model_for_Chinese_International_Education_Teaching_Courses} \\
        \hline
        Course Assessment and Student Satisfaction & Course Duration \& Difficulty, Student satisfaction, online learning efficacy during crises, Student's college experience, department reputation, course ratings, Course assessment, instructor attributes, course content, Course assessment, learner feedback, Course satisfaction, student willingness. & \cite{Experience_in_learners_review_to_determine_attribute_relation_for_course_completion}, \cite{An_Efficient_Data_Mining_Technique_for_Assessing_Satisfaction_Level_With_Online_Learning_for_Higher_Education_Students_During_the_COVID-19}, \cite{Romanov},\cite{Understanding_Students_Online_Reviews_to_Improve_College_Experience_and_Graduation_Rates_of_STEM_Programs_at_the_Largest_Post-Secondary_Institution_A_Learner-Centered_Study}, \cite{What_predicts_student_satisfaction_with_MOOCs:_A_gradient_boosting_trees_supervised_machine_learning_and_sentiment_analysis_approach}, \cite{Assessing_learners_satisfaction_in_collaborative_online_courses_through_a_big_data_approach}, \cite{Course_Satisfaction_in_Engineering_Education_Through_the_Lens_of_Student_Agency_Analytics} \\
        \hline
        Course Duration, Difficulty and Completion Rates & Curriculum quality, course duration and difficulty, completion rates, Course Duration \& Difficulty, Course dropout rate, learning habits. & \cite{Competition_Neural_Network_Model_Application_in_the_Bilingual_Teaching_Quality_Evaluation_of_Universities}, \cite{Experience_in_learners_review_to_determine_attribute_relation_for_course_completion}, \cite{Exploring_Causes_for_the_Dropout_on_Massive_Open_Online_Courses} \\
        \hline
        No mention for impact  & N/A & \cite{Research_on_the_Application_of_BP_Neural_Network_in_Vocal_Music_Teaching_Quality_Evaluation}, \cite{Design_of_online_teaching_quality_evaluation_system_for_Private_University_Research_based_on_deep_learning_algorithm}, \cite{Sentiment_analysis_of_student_feedback_using_machine_learning_and_lexicon-based_approaches}, \cite{BP_Neural_Network-Based_Evaluation_on_University_Teachers'_Teaching_Quality}, \cite{Automatic_Rules_Generation_for_Teaching_Strategies_Improvement}, \cite{Evaluation_Model_of_PE_Teaching_Quality_Based_on_Active_Learning_Support_Vector_Machine}, \cite{Evaluation_System_of_Tourism_Psychology_Teaching_Quality_Based_on_Convolutional_Neural_Network}, \cite{An_Evaluation_Model_of_Moral_Education_Teaching_Quality_in_Universities_Based_on_System_Control_Theory}, \cite{Personalized_Learning_Model_based_on_Deep_Learning_Algorithm_for_Student_Behaviour_Analytic}, \cite{The_Evaluation_and_Application_of_BP_Neural_Network_in_the_Multiple_Quality_Assurance_Paths_of_Higher_Education} \\
        \hline
    \end{tabularx}
\end{table*}

\section{SYNOPTIC SUMMARY AND SIGNIFICANCE OF RESEARCH}
\label{SYNOPTIC SUMMARY AND SIGNIFICANCE OF RESEARCH}
In our comprehensive review, we observed multiple use cases of AI application intended to enhance education, in both online and offline contexts. After an extensive study, we identified a pattern grouping these use cases together:

\begin{enumerate}
  \item Dual-Technique Sentiment Analysis for Focused Educational Feedback.
  \item Single-Technique Text Analysis for Specific Educational Enhancements.
  \item Neural Network-Based Quantitative Forecasting for Educational Outcomes.
  \item Computer Vision for Tracking Emotions.
\end{enumerate}

A detailed explanation of these groupings is available in Section \ref{RQ1_analysis}. Interestingly, sentiment analysis was found to be the most frequently utilized AI technique, with many researchers aiming to analyze feedback comments for more profound insights.

Second, several metrics were employed to measure teaching quality. We noticed these metrics can be clustered based on their utility and their area of impact. The identified clusters are:

\begin{enumerate}
  \item Feedback Analysis and Course Design.
  \item Teaching Quality and Effectiveness Papers.
  \item Course Assessment and Student Satisfaction.
  \item Course Duration, Difficulty and Completion Rates.
\end{enumerate}

We further analyzed the factors utilized within these metrics. Details of these findings are presented in Section \ref{RQ2_analysis}.

Third, we examined the impact of AI techniques used in reviewed papers, and the educational elements impacted by these techniques.Details of these findings are presented in Section \ref{RQ3_analysis}. The identified clusters of these elements are:

\begin{itemize}
  \item Feedback Analysis and Course Design.
  \item Teaching Quality and Effectiveness Papers.
  \item Course Assessment and Student Satisfaction.
  \item Course Duration, Difficulty and Completion Rates.
\end{itemize}

Lastly, not many researcher created a tool that the community can use except researches:\cite{Palaute:_An_Online_Text_Mining_Tool_for_Analyzing_Written_Student_Course_Feedback},
\cite{SUFAT_An_Analytics_Tool_for_Gaining_Insights_from_Student_Feedback_Comments},
\cite{Online_reviews_evaluation_system_for_higher_education_institution:_An_aspect_based_sentiment_analysis_tool},
\cite{Assessing_learners_satisfaction_in_collaborative_online_courses_through_a_big_data_approach}, and 
\cite{Analyzing_educational_comments_for_topics_and_sentiments:_A_text_analytics_approach}.


\section{LIMITATIONS, THREATS TO VALIDITY, AND REVIEW ASSESSMENT}
\label{LIMITATIONS, THREATS TO VALIDITY, AND REVIEW ASSESSMENT}
\subsection{Limitations}
\label{limitation}
\begin{enumerate}

\item Our study confines itself to research from the past decade, introducing a potential limitation in acquiring comprehensive knowledge. However, precedent Systematic Literature Reviews (SLRs) in this area, such as the one by \cite{AI_IN_Education_SLR}, observed that the majority of relevant studies were published between 2016 and 2019. Our analysis  of publication years resonated with these findings illustrated in figure \ref{clusterByYears}, suggesting a concentration of research activity during this period. This observation gives us confidence that our 10-year limit, while potentially excluding some earlier work, is unlikely to significantly impact the thoroughness or relevance of our findings.

\item The search was carried out utilizing a limited number of databases, namely IEEExplore, ScienceDirect, and ACM Digital Library. Some readers may argue that we should have considered additional databases for our searches, such as Google Scholar.We observed that Google Scholar often simply compiles papers from other databases, sometimes with no filter for academic quality. This leads to low-quality conference proceedings and questionable gray literature often being found through Google Scholar. Although that was our selection of database but we strongly believe that this selection helped us to insure solid academic findings 


    \item While conducting this systematic literature review, we acknowledge the potential shortcoming of not considering grey literature. It is common for SLRs to prioritize primary studies over grey literature, which could potentially lead to missing out on relevant findings. However, in order to ensure the quality and rigor of our review process, we made the decision to exclude grey literature from our analysis. While this may have limited the scope of our review, it allowed us to focus on high-quality studies and provide a more rigorous evaluation of the research in the field.

    \item To conclude, it could be contended that the decision to solely include papers written in English, as one of the inclusion criteria in this study, could be considered as highly limiting the scope of research available. Nevertheless, it is important to note that a significant portion of the existing literature in the field is written in English, hence this criterion is not peculiar or rare in the context of our field of research.


\end{enumerate}

\subsection{Threats to Validity}
This section of the paper aims to address any possible biases that may have influenced our research and the results we obtained. We will outline the measures we took to prevent the occurrence of such biases.\cite{A_NewParadigm_for_SystematicLiteratureReviews}

\begin{enumerate}
    \item \textbf{Retrieval bias} - In order to avoid any potential bias, we carried out extensive searches across three different databases. We also made sure to select only high-quality articles, and we only included works that were relevant to our research topic. although we might be biased and limited the search by last 10 years of publication.

    \item \textbf{Publication bias} - Our research was based on high-quality, peer-reviewed articles that were published in credible academic journals or distinguished conferences. We aimed to cover a range of analysis levels and experimental protocols, to ensure that we did not solely report statistically significant results.
    \item \textbf{Inclusion criteria bias} - To prevent bias, we endeavored to develop the most appropriate and inclusive criteria for our subject matter, adhering to the highest standards of our field.
    \item \textbf{Selector bias} -we thoroughly reviewed all the papers that passed our exclusion criteria to ensure the validity of our selection. Some papers were deemed obviously irrelevant to our research, and were therefore excluded. Other papers were selected based on the aforementioned inclusion criteria, and we took care to read through the methodology, discussion, and results of each selected paper to ensure that we obtained the best results that matched our findings. This rigorous process was adopted to guarantee the quality and accuracy of our research.
\end{enumerate}

\subsection{Review Assessment}
The final stage of evaluating our research findings involves an overall assessment of the quality of our work. As suggested in \cite{Procedures_for_PerformingSystematicReviews}, we have created a list of questions that can be used as a benchmark for assessing the quality of an SLR. These questions are provided below for reference.

\begin{enumerate}

\item \textbf{Were the inclusion and exclusion criteria based on reasonable and objective grounds?}
Before starting the search process, we mentioned all the criteria in the methodology. The criteria we chose are in agreement with the ones usually used in the field and are appropriate for the topic of our research. Therefore, we consider the inclusion/exclusion criteria we utilized to be rational.
\item \textbf{Was the search process comprehensive and covered all potentially relevant papers?}
We conducted our searches using trustworthy bibliographic databases and collected papers from highly-regarded academic journals and reputable conference proceedings. Based on this approach, we can state that the papers chosen for inclusion in our final record are indicative of the field.

\item \textbf{Was there a quality review of the selected papers?}
We created a measurement standard to evaluate the quality of the papers. As a matter of fact, we demonstrated that the papers included in the record possessed a high level of quality. Hence, we can affirm that the conclusions drawn from our systematic literature review were dependable and precise.
\end{enumerate}


\section{conclusion}
\label{conclusion}
In concluding our comprehensive research, we have established key findings and potential directions for future investigation in the field of AI application in education, informed by the analysis of 30 diverse papers through a systematic literature review (SLR).

Firstly, regarding the use-cases of AI in education (RQ1), four prominent approaches emerged. These included dual-technique sentiment analysis, single-technique text analysis, neural network-based quantitative forecasting, and computer vision applications. Sentiment analysis emerged as the most prevalent technique, focusing primarily on mining and interpreting feedback comments to provide more comprehensive insights, potentially leading to more focused and effective educational interventions. However, it's clear that these techniques are not mutually exclusive and can be combined to yield richer insights and more nuanced applications in education.

In terms of measuring the quality of education (RQ2), several metrics were identified, falling into clusters of Feedback Analysis and Course Design, Teaching Quality and Effectiveness, Course Assessment and Student Satisfaction, and Course Duration, Difficulty and Completion Rates. These metrics address a range of aspects in both online and offline education, suggesting that the quality of education is a multidimensional construct that requires a multifaceted approach to measure effectively. Each cluster has its own set of factors that contribute to their evaluation, indicating the depth and complexity of what constitutes 'quality' in education.

Thirdly, the impact of AI on educational quality (RQ3) was found to be profound across both online and offline formats. The same clusters as in RQ2 emerged, demonstrating that AI can be utilized effectively across these diverse aspects of educational quality. The ability of AI to analyze feedback and design courses, assess teaching quality and effectiveness, evaluate course satisfaction, and track course duration, difficulty and completion rates, underscores its transformative potential in enhancing educational quality and student outcomes.


Overall, the breadth and depth of AI application in education are vast, and its potential to enhance quality is substantial. As AI techniques continue to evolve and mature, their influence on education is expected to grow, warranting continuous study and refinement. Future research should also seek to address existing gaps and explore new AI techniques to further the understanding and application of AI in education.

% \section{Acknowledgment}
% \label{Acknowledgment}
% this research has been supported by Innopolis university
\section{Appendix}
\label{Appendix}

\subsection{Prisma checklist}
the Prisma check list can be found in the table \ref{Prisma checklist}.
% Please add the following required packages to your document preamble:
% \usepackage[table,xcdraw]{xcolor}
% If you use beamer only pass "xcolor=table" option, i.e. \documentclass[xcolor=table]{beamer}
\begin{table}[]
\begin{tabular}{|ll|l|}
\hline
\multicolumn{2}{|c|}{{\color[HTML]{000000} \textbf{Number}}} & \multicolumn{1}{c|}{{\color[HTML]{000000} \textbf{Location}}} \\ \hline
\multicolumn{2}{|l|}{\textbf{TITLE}}                         &                                                               \\ \hline
\multicolumn{2}{|l|}{1}                                      & page \ref{Abstract}                                                             \\ \hline
\multicolumn{2}{|l|}{\textbf{ABSTRACT}}                      &                                                               \\ \hline
\multicolumn{2}{|l|}{2}                                      & section \ref{Abstract}                                                             \\ \hline
\multicolumn{2}{|l|}{\textbf{INTRODUCTION}}                  &                                                               \\ \hline
\multicolumn{2}{|l|}{3}                                      & section \ref{introduction}                                                              \\ \hline
\multicolumn{2}{|l|}{4}                                      & section \ref{introduction}                                                               \\ \hline
\multicolumn{2}{|l|}{\textbf{METHODS}}                       &                                                               \\ \hline
\multicolumn{2}{|l|}{5}                                      & section \ref{Review PROTOCOL}                                                       \\ \hline
\multicolumn{2}{|l|}{6}                                      & section \ref{searchStrategy}                                           \\ \hline
\multicolumn{2}{|l|}{7}                                      & section \ref{searchStrategy}                                          \\ \hline
\multicolumn{2}{|l|}{8-15}                                   & not applicable                                                \\ \hline
\multicolumn{2}{|l|}{16a}                                    & section \ref{Review PROTOCOL}                                                        \\ \hline
\multicolumn{2}{|l|}{16b}                                    & not applicable                                                \\ \hline
\multicolumn{2}{|l|}{\textbf{RESULTS}}                       &                                                               \\ \hline
\multicolumn{2}{|l|}{17}                                     & section \ref{Results and Discussion}                                                    \\ \hline
\multicolumn{2}{|l|}{18-22}                                  & not applicable                                                \\ \hline
\multicolumn{2}{|l|}{\textbf{DISCUSSION}}                    &                                                               \\ \hline
\multicolumn{2}{|l|}{23a}                                    & sections[\ref{RQ1_analysis},\ref{RQ2_analysis},\ref{RQ3_analysis}]                                             \\ \hline
\multicolumn{2}{|l|}{23b}                                    & section \ref{Results and Discussion}                                               \\ \hline
\multicolumn{2}{|l|}{\textbf{LIMITATION}}                   &                                                               \\ \hline
\multicolumn{2}{|l|}{23c}                                    & section \ref{LIMITATIONS, THREATS TO VALIDITY, AND REVIEW ASSESSMENT}                                                 \\ \hline
\multicolumn{2}{|l|}{\textbf{CONCLUSION}}                    &                                                               \\ \hline
\multicolumn{2}{|l|}{23d}                                    & section \ref{conclusion}                                                 \\ \hline
\multicolumn{2}{|l|}{\textbf{OTHER INFORMATION}}             &                                                               \\ \hline
\multicolumn{2}{|l|}{24}                                     & not applicable                                                \\ \hline
\multicolumn{2}{|l|}{25}                                     & missing \\ \hline
\end{tabular}
\caption{Prisma checklist}
\label{Prisma checklist}
\end{table}

\subsection{Quality Assessment}
\label{Quality_Assessment}
our evaluating of the paper is out of 3 and the average quality of the papers is 2.53. all the quality of the papers presented in table \ref{tab:Quality assessment}. detaile explination of the assessment points are below: every paper will take 1 point out of the three categories below thus accumlating 3 


\begin{enumerate}
    \item
    \textbf{Q1- Does the research provide a detailed description how they used AI technology in their research?}
\begin{enumerate}
\item 1 - if the research provide full description               
\item 0 - if the research intentionally hide how they are solving the problem           
\end{enumerate}


\item 
\textbf{Q2- Does the research provide detail on the research design methodology and how the reach the results?}\begin{enumerate}
\item  1: authors provided full research
\item  0: authors didn't provide deep explanation of their methodology 
\end{enumerate}

\item 
\textbf{Q3-Does the research provide an explanation how this techniques impacted the educational institution?}
\begin{enumerate}
\item  1- if the authors provided explanation of the impact on the institute                  
\item  0 if the authors did not provided explanation of the impact on the institute          
\end{enumerate}
\end{enumerate}
\begin{table}[]
\begin{tabular}{|l|l|l|}
\hline
Paper References                                                                                                                                                                    & Rating & Missing Quality                                                 \\ \hline
\cite{Palaute:_An_Online_Text_Mining_Tool_for_Analyzing_Written_Student_Course_Feedback}                                                                                            & 3      & -                                                               \\ \hline
\cite{Experience_in_learners_review_to_determine_attribute_relation_for_course_completion}                                                                                          & 3      & -                                                               \\ \hline
\cite{Competition_Neural_Network_Model_Application_in_the_Bilingual_Teaching_Quality_Evaluation_of_Universities}                                                                    & 2      & Missing Q3                                                      \\ \hline
\cite{SUFAT_An_Analytics_Tool_for_Gaining_Insights_from_Student_Feedback_Comments}                                                                                                  & 3      & -                                                               \\ \hline
\cite{Research_on_the_Application_of_BP_Neural_Network_in_Vocal_Music_Teaching_Quality_Evaluation}                                                                                  & 2      & Missing Q3                                                      \\ \hline
\cite{An_Efficient_Data_Mining_Technique_for_Assessing_Satisfaction_Level_With_Online_Learning_for_Higher_Education_Students_During_the_COVID-19}                                   & 3      & -                                                               \\ \hline
\cite{Design_of_online_teaching_quality_evaluation_system_for_Private_University_Research_based_on_deep_learning_algorithm}                                                         & 1      & \begin{tabular}[c]{@{}l@{}}Missing Q3\\ Missing Q2\end{tabular} \\ \hline
\cite{Sentiment_analysis_of_student_feedback_using_machine_learning_and_lexicon-based_approaches}                                                                                   & 1      & \begin{tabular}[c]{@{}l@{}}Missing Q3\\ Missing Q2\end{tabular} \\ \hline
\cite{Understanding_Students_Online_Reviews_to_Improve_College_Experience_and_Graduation_Rates_of_STEM_Programs_at_the_Largest_Post-Secondary_Institution_A_Learner-Centered_Study} & 3      & -                                                               \\ \hline
\cite{Evaluating_the_quality_of_teaching_in_higher_education_institutes_using_clustering_approach_and_set_pair_analysis}                                                            & 3      & -                                                               \\ \hline
\cite{A_Study_of_MOOC_Course_Review_Topics_Mining_Based_on_LDA_Topic_Model}                                                                                                         & 3      & -                                                               \\ \hline
\cite{BP_Neural_Network-Based_Evaluation_on_University_Teachers'_Teaching_Quality}                                                                                                  & 2      & -                                                               \\ \hline
\cite{An_Improved_BP_Neural_Network-Based_Quality_Evaluation_Model_for_Chinese_International_Education_Teaching_Courses}                                                            & 2      & Missing Q2                                                      \\ \hline
\cite{Automatic_Rules_Generation_for_Teaching_Strategies_Improvement}                                                                                                               & 2      & Missing Q2                                                      \\ \hline
\cite{Evaluation_Model_of_PE_Teaching_Quality_Based_on_Active_Learning_Support_Vector_Machine}                                                                                      & 2      & Missing Q2                                                      \\ \hline
\cite{Evaluation_System_of_Tourism_Psychology_Teaching_Quality_Based_on_Convolutional_Neural_Network}                                                                               & 2      & Missing Q2                                                      \\ \hline
\cite{What_predicts_student_satisfaction_with_MOOCs:_A_gradient_boosting_trees_supervised_machine_learning_and_sentiment_analysis_approach}                                         & 3      & -                                                               \\ \hline
\cite{Assessing_learners_satisfaction_in_collaborative_online_courses_through_a_big_data_approach}                                                                                  & 3      & -                                                               \\ \hline
\cite{Analyzing_educational_comments_for_topics_and_sentiments:_A_text_analytics_approach}                                                                                          & 3      & -                                                               \\ \hline
\cite{Research_on_Project-Based_Teaching_Quality_Evaluation_Based_on_Fuzzy_Cluster_Analysis}                                                                                        & 3      & -                                                               \\ \hline
\cite{Improving_MOOC_quality_using_learning_analytics_and_tools}                                                                                                                    & 3      & -                                                               \\ \hline
\cite{Visualizing_Student_Opinion_Through_Text_Analysis}                                                                                                                            & 3      & -                                                               \\ \hline
\cite{Course_Satisfaction_in_Engineering_Education_Through_the_Lens_of_Student_Agency_Analytics}                                                                                    & 3      & -                                                               \\ \hline
\cite{An_Evaluation_Model_of_Moral_Education_Teaching_Quality_in_Universities_Based_on_System_Control_Theory}                                                                       & 2      & Missing Q3                                                      \\ \hline
\cite{Research_on_Evaluation_of_Online_Teaching_Effect_based_on_Deep_Learning_Technology}                                                                                           & 3      & -                                                               \\ \hline
\cite{Exploring_Causes_for_the_Dropout_on_Massive_Open_Online_Courses}                                                                                                              & 3      & -                                                               \\ \hline
\cite{Personalized_Learning_Model_based_on_Deep_Learning_Algorithm_for_Student_Behaviour_Analytic}                                                                                  & 2      & Missing Q2                                                      \\ \hline
\cite{Research_on_the_Application_of_BP_Neural_Network_in_the_Evaluation_of_College_Teaching_Quality}                                                                               & 3      & -                                                               \\ \hline
\cite{The_Evaluation_and_Application_of_BP_Neural_Network_in_the_Multiple_Quality_Assurance_Paths_of_Higher_Education}                                                              & 2      & Missing Q3                                                      \\ \hline
\cite{Online_reviews_evaluation_system_for_higher_education_institution:_An_aspect_based_sentiment_analysis_tool}                                                                   & 3      & -                                                               \\ \hline
\end{tabular}
\caption{Quality assessment}
\label{tab:Quality assessment}
\end{table}



\bibliographystyle{plain}
\bibliography{references}


% ---- Bibliography ----
%
% \begin{thebibliography}{99}

% \bibitem{An_Evaluation_Model_of_Moral_Education_Teaching_Quality_in_Universities_Based_on_System_Control_Theory} 
% Ming, Zheng. "An Evaluation Model of Moral Education Teaching Quality in Universities Based on System Control Theory." 2021 13th International Conference on Measuring Technology and Mechatronics Automation (ICMTMA), 2021, pp. 535-538. \url{https://doi.org/10.1109/ICMTMA52658.2021.00123}.

% \bibitem{Procedures_for_PerformingSystematicReviews} 
% Kitchenham, Barbara. "Procedures for Performing Systematic Reviews." Keele University, UK, 2004, Vol. 33.

% \bibitem{Course_Satisfaction_in_Engineering_Education_Through_the_Lens_of_Student_Agency_Analytics} 
% Heilala, Ville, et al. "Course Satisfaction in Engineering Education Through the Lens of Student Agency Analytics." 2020 IEEE Frontiers in Education Conference (FIE), 2020, pp. 1-9. \url{https://doi.org/10.1109/FIE44824.2020.9274141}.

% \bibitem{Assessing_learners_satisfaction_in_collaborative_online_courses_through_a_big_data_approach} 
% Elia, Gianluca, et al. "Assessing learners’ satisfaction in collaborative online courses through a big data approach." Computers in Human Behavior, vol. 92, 2019, pp. 589-599. \url{https://doi.org/10.1016/j.chb.2018.04.033}.

% \bibitem{What_predicts_student_satisfaction_with_MOOCs:_A_gradient_boosting_trees_supervised_machine_learning_and_sentiment_analysis_approach} 
% Hew, Khe Foon, et al. "What predicts student satisfaction with MOOCs: A gradient boosting trees supervised machine learning and sentiment analysis approach." Computers \& Education, vol. 145, 2020, p. 103724. \url{https://doi.org/10.1016/j.compedu.2019.103724}.

% \bibitem{Automatic_Rules_Generation_for_Teaching_Strategies_Improvement} 
% Zhan, Jing, et al. "Automatic Rules Generation for Teaching Strategies Improvement." Proceedings of the 4th International Conference on Computer Science and Application Engineering, 2020, Article 97. \url{https://doi.org/10.1145/3424978.3425079}.

% \bibitem{Evaluation_System_of_Tourism_Psychology_Teaching_Quality_Based_on_Convolutional_Neural_Network} 
% Chen, Peilin. "Evaluation System of Tourism Psychology Teaching Quality Based on Convolutional Neural Network." Proceedings of the 2021 4th International Conference on Information Systems and Computer Aided Education, 2021, pp. 964-967. \url{https://doi.org/10.1145/3482632.3483062}.

% \bibitem{Evaluation_Model_of_PE_Teaching_Quality_Based_on_Active_Learning_Support_Vector_Machine} 
% Zhang, Fan, et al. "Evaluation Model of PE Teaching Quality Based on Active Learning Support Vector Machine." Proceedings of the 5th International Conference on Frontiers of Educational Technologies, 2019, pp. 101-104. \url{https://doi.org/10.1145/3338188.3338202}.

% \bibitem{BP_Neural_Network-Based_Evaluation_on_University_Teachers'_Teaching_Quality} 
% Guo, Chujian, and Yicai Liu. "BP Neural Network-Based Evaluation on University Teachers' Teaching Quality." Proceedings of the 2020 3rd International Conference on E-Business, Information Management and Computer Science, 2021, pp. 285-289. \url{https://doi.org/10.1145/3453187.3453348}.

% \bibitem{Exploring_Causes_for_the_Dropout_on_Massive_Open_Online_Courses} 
% Niu, Zhiyao, et al. "Exploring Causes for the Dropout on Massive Open Online Courses." Proceedings of ACM Turing Celebration Conference - China, 2018, pp. 47-52. \url{https://doi.org/10.1145/3210713.3210727}.

% \bibitem{Research_on_the_Application_of_BP_Neural_Network_in_the_Evaluation_of_College_Teaching_Quality} 
% Shen, Xin, and Tao Bian. "Research on the Application of BP Neural Network in the Evaluation of College Teaching Quality." Proceedings of the 2022 6th International Conference on E-Education, E-Business and E-Technology, 2022, pp. 73-77. \url{https://doi.org/10.1145/3549843.3549854}.

% \bibitem{A_Structural_Examination_and_Exploration_of_Chinese_University_Students’_Satisfaction_with_MOOCs} 
% Meng, Ling, et al. "A Structural Examination and Exploration of Chinese University Students’ Satisfaction with MOOCs." Proceedings of the 7th International Conference on Distance Education and Learning, 2022, pp. 190-196. \url{https://doi.org/10.1145/3543321.3543353}.

% \bibitem{AI_IN_Education_SLR} 
% Tahiru, Fati. "AI in Education: A Systematic Literature Review." Journal of Cases on Information Technology, vol. 23, 2021, pp. 1-20. \url{https://doi.org/10.4018/JCIT.2021010101}.

% \bibitem{Technologies_for_GQM_Based_Metrics_Recommender_Systems_A_Systematic_Literature_Review} 
% Farina, Mirko, et al. "Technologies for GQM-Based Metrics Recommender Systems: A Systematic Literature Review." IEEE Access, vol. 10, 2022, pp. 23098-23111. \url{https://doi.org/10.1109/ACCESS.2022.3152397}.

% \end{thebibliography}
\end{document}
